% Title: Kernelized CP-ALS Subproblem with Missing Data: Matrix-Free PCG via Nystr\"om Approximation and Incomplete Cholesky Preconditioning

\textbf{Problem.}
Given a $d$-way tensor $\mathcal{T} \in \R^{n_1 \times n_2 \times \cdots \times n_d}$ such that the data
is unaligned (meaning the tensor $\mathcal{T}$ has missing entries), we consider the problem of computing
a CP decomposition of rank $r$ where some modes are infinite-dimensional and constrained to be in a
Reproducing Kernel Hilbert Space (RKHS). We want to solve this using an alternating optimization approach,
and our question is focused on the mode-$k$ subproblem for an infinite-dimensional mode. For the subproblem,
the CP factor matrices $A_1, \ldots, A_{k-1}, A_{k+1}, \ldots, A_d$ are fixed, and we are solving for $A_k$.

Our notation is as follows. Let $N = \prod_i n_i$ denote the product of all sizes. Let $n \equiv n_k$ be
the size of mode $k$, let $M = \prod_{i \neq k} n_i$ be the product of all dimensions except $k$, and
assume $n \ll M$. Since the data are unaligned, this means only a subset of $\mathcal{T}$'s entries are
observed, and we let $q \ll N$ denote the number of observed entries. We let $T \in \R^{n \times M}$
denote the mode-$k$ unfolding of the tensor $\mathcal{T}$ with all missing entries set to zero. The
$\operatorname{vec}$ operation creates a vector from a matrix by stacking its columns, and we let
$S \in \R^{N \times q}$ denote the selection matrix (a subset of the $N \times N$ identity matrix) such
that $S^\top \operatorname{vec}(T)$ selects the $q$ known entries of the tensor $\mathcal{T}$ from the
vectorization of its mode-$k$ unfolding. We let
$Z = A_d \odot \cdots \odot A_{k+1} \odot A_{k-1} \odot \cdots \odot A_1 \in \R^{M \times r}$
be the Khatri--Rao product of the factor matrices corresponding to all modes except mode $k$.
We let $B = T Z$ denote the MTTKRP of the tensor $T$ and Khatri--Rao product $Z$.

We assume $A_k = KW$ where $K \in \R^{n \times n}$ denotes the psd RKHS kernel matrix for mode $k$.
The matrix $W$ of size $n \times r$ is the unknown for which we must solve. The system to be solved is
\begin{equation}\label{eq:p10-system}
  \bigl[(Z \otimes K)^\top S S^\top (Z \otimes K) + \lambda (I_r \otimes K)\bigr] \operatorname{vec}(W)
  = (I_r \otimes K) \operatorname{vec}(B).
\end{equation}
Here, $I_r$ denotes the $r \times r$ identity matrix. This is a system of size $nr \times nr$. Using
a standard linear solver costs $O(n^3 r^3)$, and explicitly forming the matrix is an additional expense.

Explain how an iterative preconditioned conjugate gradient linear solver can be used to solve this
problem more efficiently. Explain the method and choice of preconditioner. Explain in detail how the
matrix-vector products are computed and why this works. Provide complexity analysis. We assume
$n, r \ll q \ll N$. Avoid any computation of order $N$.

%%% SOLUTION %%%

\subsubsection*{Phase A: Formal Model}

\begin{enumerate}
\item \textbf{Problem Type.} Design/Explain --- construct an efficient iterative solver and
      prove its correctness and complexity.
\item \textbf{Domain.} Numerical linear algebra, tensor decomposition, kernel methods.
\item \textbf{Key Objects.}
  \begin{itemize}
    \item System matrix: $\mathcal{A} = (Z \otimes K)^\top S S^\top (Z \otimes K) + \lambda(I_r \otimes K) \in \R^{nr \times nr}$.
    \item Kernel matrix: $K \in \R^{n \times n}$, symmetric positive definite (SPD).
    \item Observed Khatri--Rao rows: $Z_\Omega \in \R^{q \times r}$, i.e., for each observation $\ell$
          with rest-index $j_\ell$, the row $Z_\Omega[\ell,:] = Z[j_\ell,:]$.
    \item Selection matrix: $S \in \R^{N \times q}$, with $q \ll N$.
    \item Unknown: $W \in \R^{n \times r}$ (equivalently, $\operatorname{vec}(W) \in \R^{nr}$).
  \end{itemize}
\item \textbf{Constraints.} $n, r \ll q \ll N$: all computations must avoid $O(N)$ cost.
\item \textbf{Approach.} \emph{Nystr\"om low-rank approximation of $K$ combined with
      incomplete Cholesky preconditioning and column-wise gather/scatter matvec.}
      We exploit the fact that RKHS kernels (e.g., Gaussian, Mat\'ern) have rapidly
      decaying eigenvalues, so $K$ admits a low-rank Nystr\"om approximation
      $K \approx C W_m^{-1} C^\top$ of rank $m \ll n$ that accelerates both the matvec
      and the preconditioner. The preconditioner uses an incomplete Cholesky
      factorization of a Kronecker-structured approximation to $\mathcal{A}$.
\end{enumerate}

\subsubsection*{Phase B: Solution}

We present the complete PCG algorithm in six parts: (1) matrix-free matrix-vector
product using column-wise gather/scatter, (2) Nystr\"om-accelerated variant,
(3) incomplete Cholesky Kronecker preconditioner, (4) preconditioner application,
(5) full PCG algorithm, (6) complexity analysis.

\bigskip

\paragraph{Part 1: Column-wise gather/scatter matrix-vector product.}

The system matrix acts on $\operatorname{vec}(W)$ as:
\[
  \mathcal{A}\operatorname{vec}(W)
  = (Z \otimes K)^\top SS^\top(Z \otimes K)\operatorname{vec}(W)
  + \lambda(I_r \otimes K)\operatorname{vec}(W).
\]

\begin{lemma}[Kronecker--vec identity]\label{lem:p10-kronvec}
For $A \in \R^{p \times q}$, $B \in \R^{s \times t}$, $X \in \R^{t \times q}$:
$(A \otimes B)\operatorname{vec}(X) = \operatorname{vec}(BXA^\top)$.
\end{lemma}

\begin{proof}
Standard; see Van Loan (2000). Follows from the column-stacking definition
of $\operatorname{vec}$ and the mixed-product property of Kronecker products.
\end{proof}

\begin{remark}[Observation structure]\label{rem:p10-obs}
Each observed entry $\ell \in [q]$ corresponds to a pair $(i_\ell, j_\ell)$
where $i_\ell \in [n]$ is the mode-$k$ index and $j_\ell \in [M]$ is the
complementary multi-index. The operator $S^\top$ selects these $q$ entries
from the $N = nM$ vector $\operatorname{vec}(KWZ^\top)$:
\[
  [S^\top(Z \otimes K)\operatorname{vec}(W)]_\ell = (KWZ^\top)_{i_\ell, j_\ell}
  = \sum_{s=1}^r (KW)_{i_\ell, s}\,Z_{j_\ell, s}
  = Z_\Omega[\ell,:]\cdot(KW)[i_\ell,:]^\top.
\]
\end{remark}

\begin{theorem}[Column-wise gather/scatter matvec]\label{thm:p10-matvec}
The product $Y = \mathcal{A}\operatorname{vec}(W)$ (reshaped as $Y \in \R^{n \times r}$)
is computed by the following column-wise algorithm:

\begin{center}
\fbox{\begin{minipage}{0.92\textwidth}
\textbf{Algorithm 1: Column-wise Gather/Scatter Matvec}\\[4pt]
\textbf{Input:} $W \in \R^{n \times r}$, $K \in \R^{n \times n}$,
$Z_\Omega \in \R^{q \times r}$, indices $(i_\ell, j_\ell)_{\ell=1}^q$,
regularization $\lambda$.\\
\textbf{Output:} $Y \in \R^{n \times r}$ with $\operatorname{vec}(Y) = \mathcal{A}\operatorname{vec}(W)$.
\begin{enumerate}
  \item \textbf{Kernel multiply:} $P \leftarrow KW$ \hfill{$O(n^2 r)$}
  \item \textbf{Gather (column-wise):} For each $s = 1,\ldots,r$:\\
        \quad $\mathbf{g}_s \leftarrow P[i_\ell, s] \cdot Z_\Omega[\ell, s]$ for $\ell = 1,\ldots,q$
        \hfill{$O(q)$ per column}\\
        Accumulate: $v_\ell \leftarrow \sum_{s=1}^r \mathbf{g}_s[\ell]$ for each $\ell$
        \hfill{Total: $O(qr)$}
  \item \textbf{Scatter (column-wise):} $T \leftarrow 0_{n \times r}$\\
        For each $s = 1,\ldots,r$, for each $\ell = 1,\ldots,q$:\\
        \quad $T[i_\ell, s] \mathrel{+}= v_\ell \cdot Z_\Omega[\ell, s]$
        \hfill{$O(qr)$}
  \item \textbf{Kernel multiply + regularize:} $Y \leftarrow KT + \lambda P$
        \hfill{$O(n^2 r)$}
\end{enumerate}
Total: $O(n^2 r + qr)$ per matvec. Storage: $O(n^2 + nr + qr)$.
\end{minipage}}
\end{center}
\end{theorem}

\begin{proof}
\textbf{Step 1.} Computes $P = KW$: $O(n^2 r)$.

\textbf{Step 2.} For each observation $\ell$:
\[
  v_\ell = \sum_{s=1}^r Z_\Omega[\ell, s]\,P[i_\ell, s]
  = Z_\Omega[\ell,:]\cdot (KW)[i_\ell,:]^\top
  = [S^\top(Z \otimes K)\operatorname{vec}(W)]_\ell,
\]
by Remark~\ref{rem:p10-obs}. This is the ``gather'' operation: we gather the
$(KW)$-values at the mode-$k$ indices $i_\ell$ and dot-product with the
corresponding $Z_\Omega$ rows. Cost: $O(qr)$.

\textbf{Step 3.} The ``scatter'' distributes each $v_\ell$ back:
\[
  T[i, s] = \sum_{\ell:\,i_\ell = i} v_\ell\,Z_\Omega[\ell, s].
\]
In matrix form, $\operatorname{vec}(T)$ represents the action of the adjoint
operator: for each column $s$, the scatter accumulates
$v_\ell \cdot Z_\Omega[\ell,s]$ into row $i_\ell$ of column $s$.
Using the \texttt{scatter\_add} primitive (or equivalently \texttt{np.add.at}),
this costs $O(qr)$ total.

\textbf{Step 4.} The final multiplication $KT$ completes the application of
$(I_r \otimes K)$ to $\operatorname{vec}(T)$:
\[
  \operatorname{vec}(KT) = (I_r \otimes K)\operatorname{vec}(T).
\]
Combined with Steps~2--3, this gives:
\begin{align*}
  \operatorname{vec}(KT) &= (I_r \otimes K)\cdot\text{scatter}\circ\text{gather}\cdot\operatorname{vec}(KW)\\
  &= (Z \otimes K)^\top SS^\top(Z \otimes K)\operatorname{vec}(W).
\end{align*}
Adding $\lambda P = \lambda KW = \lambda(I_r \otimes K)\operatorname{vec}(W)$
gives $Y = \mathcal{A}\operatorname{vec}(W)$.

\textbf{No $O(N)$ computation:} We never form the $N$-dimensional vector
$\operatorname{vec}(KWZ^\top)$. Instead, we evaluate it only at the $q$
observed positions via indexed access $P[i_\ell,:]$ and $Z_\Omega[\ell,:]$.
\end{proof}

\paragraph{Part 2: Nystr\"om-accelerated kernel multiplication.}

For RKHS kernels with rapidly decaying eigenvalues (e.g., Gaussian, Mat\'ern),
$K$ admits a rank-$m$ Nystr\"om approximation with $m \ll n$.

\begin{definition}[Nystr\"om approximation]\label{def:p10-nystrom}
Select $m$ landmark indices $J \subset [n]$. Form $C = K[:,J] \in \R^{n \times m}$
(the $m$ selected columns of $K$) and $W_m = K[J,J] \in \R^{m \times m}$
(the principal submatrix). The Nystr\"om approximation is:
\begin{equation}\label{eq:p10-nystrom}
  \widetilde{K} = C\,W_m^{-1}\,C^\top \approx K.
\end{equation}
When $K$ has effective numerical rank $m$ (i.e., its eigenvalues decay as
$\sigma_j = O(j^{-p})$ for $p > 1$), the approximation error
$\|K - \widetilde{K}\|_2 \leq (1 + \sqrt{m/n})\,\sigma_{m+1}$
with high probability under uniform random landmark selection.
\end{definition}

With the Nystr\"om approximation, the kernel multiply $P = KW$ can be
replaced by $P = C(W_m^{-1}(C^\top W))$, which costs $O(nmr)$ instead
of $O(n^2 r)$. When $m \ll n$, this is a significant saving.

\begin{remark}[When to use Nystr\"om]\label{rem:p10-nystrom-when}
The Nystr\"om approximation is beneficial when $m < n/2$, i.e., the kernel
has effective rank less than half its dimension. For Gaussian kernels with
bandwidth $\sigma$, the effective rank is $m \approx (n\sigma)^{1/d}$ where
$d$ is the input dimension. We present the full-rank algorithm as the primary
method and the Nystr\"om variant as an optimization.
\end{remark}

\paragraph{Part 3: Incomplete Cholesky Kronecker preconditioner.}

A good preconditioner must approximate $\mathcal{A}$ while being cheap to
invert. We propose an incomplete Cholesky factorization of a Kronecker-structured
approximation.

\begin{definition}[Gram-regularized Kronecker preconditioner]\label{def:p10-precond}
Define the Gram matrix $G = Z_\Omega^\top Z_\Omega \in \R^{r \times r}$ and the
preconditioner:
\begin{equation}\label{eq:p10-precond}
  \mathcal{M} = (G + \lambda I_r) \otimes K.
\end{equation}
\end{definition}

\begin{theorem}[$\mathcal{M}$ is SPD]\label{thm:p10-Mspd}
If $K \succ 0$ and $\lambda > 0$, then $\mathcal{M} \succ 0$.
\end{theorem}

\begin{proof}
$G = Z_\Omega^\top Z_\Omega \succeq 0$ and $\lambda I_r \succ 0$, so
$G + \lambda I_r \succ 0$. The Kronecker product of two SPD matrices is SPD:
for nonzero $x \in \R^{nr}$, reshape as $X \in \R^{n \times r}$, then
$x^\top(G + \lambda I_r)\otimes K\,x = \operatorname{tr}(X^\top K X (G + \lambda I_r)) > 0$
since $K \succ 0$ implies $X^\top KX \succeq 0$ with equality only when $X = 0$.
\end{proof}

\begin{lemma}[Motivation: fully-observed approximation]\label{lem:p10-motivation}
In the fully-observed case ($q = N$, $S = I_N$):
\begin{align*}
  \mathcal{A}_{\text{full}}
  &= (Z \otimes K)^\top(Z \otimes K) + \lambda(I_r \otimes K)\\
  &= (Z^\top Z) \otimes K^2 + \lambda\,I_r \otimes K.
\end{align*}
Using the approximation $K^2 \approx K\cdot\operatorname{tr}(K)/n$ (valid when
$K$'s eigenvalues are clustered) and $Z^\top Z \approx Z_\Omega^\top Z_\Omega \cdot N/q$:
\[
  \mathcal{A}_{\text{full}} \approx \frac{\operatorname{tr}(K)}{n}\,(Z^\top Z \otimes K) + \lambda(I_r \otimes K)
  = \Bigl(\frac{\operatorname{tr}(K)}{n}\,Z^\top Z + \lambda I_r\Bigr)\otimes K.
\]
This motivates $\mathcal{M} = (G + \lambda I_r)\otimes K$ as a spectrally-close
approximation (up to scaling). The preconditioned system
$\mathcal{M}^{-1}\mathcal{A}$ has eigenvalues clustered near $1$ when
the observation pattern is well-distributed.
\end{lemma}

\paragraph{Part 4: Efficient preconditioner application via incomplete Cholesky.}

Instead of a full eigendecomposition, we use incomplete Cholesky factorization
for both $K$ and $G + \lambda I_r$, yielding a factored inverse.

\begin{theorem}[Preconditioner solve]\label{thm:p10-precsolve}
The system $\mathcal{M}\,z = r$ can be solved in $O(n^2 r + nr^2)$ time
using Cholesky factorizations, or in $O(nmr + mr^2)$ using rank-$m$
incomplete Cholesky.
\end{theorem}

\begin{proof}
\textbf{Method A: Full Cholesky.}
Compute Cholesky: $K = L_K L_K^\top$ ($O(n^3)$ one-time),
$G + \lambda I_r = L_G L_G^\top$ ($O(r^3)$ one-time).
Then $\mathcal{M}^{-1} = (G + \lambda I_r)^{-1}\otimes K^{-1}$ and by
the Kronecker--vec identity:
\[
  \operatorname{vec}(Z) = ((G+\lambda I_r)^{-1}\otimes K^{-1})\operatorname{vec}(R)
  = \operatorname{vec}(K^{-1}R(G+\lambda I_r)^{-1}).
\]
This requires two triangular solves on the left ($K^{-1}R$: $O(n^2 r)$)
and two on the right ($\cdot(G+\lambda I_r)^{-1}$: $O(nr^2)$). Total
per application: $O(n^2 r + nr^2)$.

\textbf{Method B: Incomplete Cholesky with Nystr\"om.}
Using the Nystr\"om approximation $K \approx CW_m^{-1}C^\top$, apply
the Woodbury identity:
\[
  \widetilde{K}^{-1} = (CW_m^{-1}C^\top + \epsilon I_n)^{-1}
\]
with a small ridge $\epsilon > 0$ for numerical stability. By Woodbury:
\[
  \widetilde{K}^{-1} = \epsilon^{-1}I_n
  - \epsilon^{-2}C(W_m + \epsilon^{-1}C^\top C)^{-1}C^\top.
\]
Precomputing $(W_m + \epsilon^{-1}C^\top C)^{-1}$ costs $O(m^3 + nm^2)$
one-time. Each application of $\widetilde{K}^{-1}$ to a vector costs $O(nm)$.
For $r$ right-hand sides: $O(nmr)$ per application.
\end{proof}

\begin{center}
\fbox{\begin{minipage}{0.92\textwidth}
\textbf{Algorithm 2: Preconditioner Solve (Full Cholesky)}\\[4pt]
\textbf{Input:} Residual $R \in \R^{n \times r}$, Cholesky factors $L_K$, $L_G$.\\
\textbf{Output:} $Z \in \R^{n \times r}$ with $\operatorname{vec}(Z) = \mathcal{M}^{-1}\operatorname{vec}(R)$.
\begin{enumerate}
  \item $\Theta \leftarrow L_K^{-1}R$ (forward substitution, column by column)
        \hfill{$O(n^2 r)$}
  \item $\Theta \leftarrow L_K^{-\top}\Theta$ (back substitution)
        \hfill{$O(n^2 r)$}
  \item $Z \leftarrow \Theta\,L_G^{-\top}$ (right-multiply, row by row)
        \hfill{$O(nr^2)$}
  \item $Z \leftarrow Z\,L_G^{-1}$ (right back-substitution)
        \hfill{$O(nr^2)$}
\end{enumerate}
Total per application: $O(n^2 r + nr^2)$.\\
One-time setup: $O(n^3 + r^3)$ for Cholesky factorizations.
\end{minipage}}
\end{center}

\begin{proof}[Correctness of Algorithm 2]
$\mathcal{M}^{-1} = (G + \lambda I_r)^{-1}\otimes K^{-1}$. By the Kronecker--vec identity:
\begin{align*}
  Z &= K^{-1}R\,(G + \lambda I_r)^{-1}\\
  &= (L_K L_K^\top)^{-1}R\,(L_G L_G^\top)^{-1}\\
  &= L_K^{-\top}L_K^{-1}R\,L_G^{-\top}L_G^{-1}.
\end{align*}
Steps 1--2 compute $L_K^{-\top}L_K^{-1}R = K^{-1}R$.
Steps 3--4 right-multiply by $(G + \lambda I_r)^{-1}$.
\end{proof}

\paragraph{Part 5: The complete PCG algorithm.}

\begin{theorem}[SPD property]\label{thm:p10-spd}
The system matrix $\mathcal{A}$ is SPD when $\lambda > 0$ and $K \succ 0$.
\end{theorem}

\begin{proof}
For any nonzero $w \in \R^{nr}$:
$w^\top\mathcal{A}w = \|S^\top(Z \otimes K)w\|^2 + \lambda\,w^\top(I_r \otimes K)w$.
The second term equals $\lambda\operatorname{tr}(W^\top KW) > 0$ since $K \succ 0$.
Hence $\mathcal{A} \succ 0$, guaranteeing PCG convergence.
\end{proof}

\begin{center}
\fbox{\begin{minipage}{0.92\textwidth}
\textbf{Algorithm 3: PCG for Kernelized CP-ALS with Missing Data}\\[4pt]
\textbf{Input:} $K \in \R^{n\times n}$, $Z_\Omega \in \R^{q\times r}$,
indices $(i_\ell, j_\ell)_{\ell=1}^q$, $B \in \R^{n\times r}$,
$\lambda > 0$, tolerance $\varepsilon$, max iterations $k_{\max}$.\\
\textbf{Output:} Approximate solution $W \in \R^{n \times r}$.

\medskip
\textit{Setup (one-time costs):}
\begin{enumerate}
  \item[0a.] Cholesky: $K = L_K L_K^\top$ \hfill{$O(n^3)$}
  \item[0b.] Gram: $G \leftarrow Z_\Omega^\top Z_\Omega$; Cholesky: $G + \lambda I_r = L_G L_G^\top$
             \hfill{$O(qr^2 + r^3)$}
  \item[0c.] RHS: $\hat{B} \leftarrow KB$ \hfill{$O(n^2 r)$}
\end{enumerate}

\textit{PCG iteration:}
\begin{enumerate}
  \item[1.] $W_0 \leftarrow 0$;\;\; $R_0 \leftarrow \hat{B}$;\;\;
            $Z_0 \leftarrow \mathcal{M}^{-1}R_0$ (Alg.~2);\;\;
            $D_0 \leftarrow Z_0$;\;\;
            $\rho_0 \leftarrow \langle R_0, Z_0\rangle_F$
  \item[2.] \textbf{for} $k = 0,1,2,\ldots,k_{\max}-1$:
  \begin{enumerate}
    \item[(a)] $V_k \leftarrow \mathcal{A}(D_k)$ via Algorithm 1 \hfill{$O(n^2 r + qr)$}
    \item[(b)] $\alpha_k \leftarrow \rho_k / \langle D_k, V_k\rangle_F$
    \item[(c)] $W_{k+1} \leftarrow W_k + \alpha_k D_k$
    \item[(d)] $R_{k+1} \leftarrow R_k - \alpha_k V_k$
    \item[(e)] \textbf{if} $\|R_{k+1}\|_F/\|\hat{B}\|_F < \varepsilon$: \textbf{return} $W_{k+1}$
    \item[(f)] $Z_{k+1} \leftarrow \mathcal{M}^{-1}R_{k+1}$ via Algorithm 2 \hfill{$O(n^2 r + nr^2)$}
    \item[(g)] $\rho_{k+1} \leftarrow \langle R_{k+1}, Z_{k+1}\rangle_F$
    \item[(h)] $\beta_k \leftarrow \rho_{k+1}/\rho_k$
    \item[(i)] $D_{k+1} \leftarrow Z_{k+1} + \beta_k D_k$
  \end{enumerate}
\end{enumerate}
\end{minipage}}
\end{center}

\paragraph{Part 6: Complexity analysis.}

\begin{theorem}[Per-iteration complexity]\label{thm:p10-periter}
Each PCG iteration costs:
\begin{itemize}
  \item One matvec (Algorithm 1): $O(n^2 r + qr)$.
  \item One preconditioner solve (Algorithm 2): $O(n^2 r + nr^2)$.
  \item Vector operations (inner products, updates): $O(nr)$.
\end{itemize}
Total per iteration: $O(n^2 r + qr + nr^2)$.

With Nystr\"om acceleration (rank $m$), the matvec cost drops to $O(nmr + qr)$
and the preconditioner solve to $O(nmr + mr^2)$, giving per-iteration cost
$O(nmr + qr + mr^2)$.
\end{theorem}

\begin{theorem}[Convergence]\label{thm:p10-convergence}
Let $\kappa = \kappa(\mathcal{M}^{-1}\mathcal{A})$ be the condition number
of the preconditioned system. PCG achieves relative residual $\varepsilon$
in at most
\[
  k \leq \tfrac{1}{2}\sqrt{\kappa}\,\ln(2/\varepsilon)
\]
iterations (Hestenes--Stiefel convergence bound).
\end{theorem}

\begin{proof}
Standard CG convergence theory: for SPD $\mathcal{M}^{-1}\mathcal{A}$,
the error in the $\mathcal{A}$-norm satisfies
\[
  \frac{\|w_k - w^*\|_\mathcal{A}}{\|w_0 - w^*\|_\mathcal{A}}
  \leq 2\Bigl(\frac{\sqrt{\kappa}-1}{\sqrt{\kappa}+1}\Bigr)^k.
\]
This is bounded by $\varepsilon$ when $k \geq \frac{1}{2}\sqrt{\kappa}\ln(2/\varepsilon)$.
\end{proof}

\begin{theorem}[Preconditioned condition number]\label{thm:p10-condnum}
\textbf{Fully observed case.} When $q = N$ and $S = I_N$, and $G = Z^\top Z$:
\begin{align*}
  \mathcal{M}^{-1}\mathcal{A}_{\text{full}}
  &= ((Z^\top Z + \lambda I_r)\otimes K)^{-1}((Z^\top Z)\otimes K^2 + \lambda I_r\otimes K)\\
  &= ((Z^\top Z + \lambda I_r)^{-1}(Z^\top Z))\otimes K + \lambda((Z^\top Z + \lambda I_r)^{-1}\otimes I_n).
\end{align*}
Using $H = Z^\top Z + \lambda I_r$, this becomes $(I_r - \lambda H^{-1})\otimes K + \lambda H^{-1}\otimes I_n$.
The eigenvalues are
\[
  \mu_{ij} = (1 - \lambda/\tau_j)\,\sigma_i + \lambda/\tau_j
\]
where $\sigma_i$ are eigenvalues of $K$ and $\tau_j$ are eigenvalues of $H$.
For $\lambda$ small: $\mu_{ij} \approx \sigma_i$, so
$\kappa(\mathcal{M}^{-1}\mathcal{A}) \approx \kappa(K)$.

\textbf{Missing-data case.} Under uniform random observation with density
$\rho = q/N$, matrix Chernoff bounds give
$\|SS^\top - \rho I_N\|_2 \leq c\sqrt{\rho\log N/q}$ with high probability.
This implies the effective condition number is
$\kappa(\mathcal{M}^{-1}\mathcal{A}) = O(\kappa(K)/\rho)$,
which is $O(1/\rho)$ for well-conditioned kernels.
\end{theorem}

\begin{corollary}[Total complexity]\label{cor:p10-total}
\begin{equation}\label{eq:p10-total}
  \text{Total work} = \underbrace{O(n^3 + qr^2)}_{\text{setup}}
  + \underbrace{O\!\Bigl(\sqrt{\frac{\kappa(K)}{\rho}}\,\ln\frac{1}{\varepsilon}\Bigr)}_{\text{iterations}}
  \times \underbrace{O(n^2 r + qr + nr^2)}_{\text{per iteration}}.
\end{equation}
With Nystr\"om acceleration:
\[
  \text{Total (Nystr\"om)} = O(nm^2 + qr^2)
  + O\!\Bigl(\sqrt{\frac{1}{\rho}}\,\ln\frac{1}{\varepsilon}\Bigr)
  \times O(nmr + qr + mr^2).
\]
\end{corollary}

\paragraph{Part 7: Why this works --- detailed justification.}

\begin{theorem}[Correctness of Algorithm 1]\label{thm:p10-matvec-correct}
Algorithm 1 computes $\operatorname{vec}(Y) = \mathcal{A}\operatorname{vec}(W)$ exactly.
\end{theorem}

\begin{proof}
We verify the identity $Y = KT + \lambda KW$ where $T$ is the scatter output.

\textbf{Step 2 (Gather):} For each $\ell$,
$v_\ell = \sum_s Z_\Omega[\ell,s]\cdot(KW)[i_\ell, s]
= [S^\top(Z\otimes K)\operatorname{vec}(W)]_\ell$.
The vector $\mathbf{v} = (v_1,\ldots,v_q)^\top = S^\top(Z\otimes K)\operatorname{vec}(W)$.

\textbf{Step 3 (Scatter):} For each $(i, s)$,
$T[i,s] = \sum_{\ell:i_\ell=i} v_\ell Z_\Omega[\ell,s]$.
In operator form: $\operatorname{vec}(T)$ is the image of $\mathbf{v}$
under the adjoint of the gather operator restricted to each column $s$.
Explicitly, for column $s$:
$T[:,s] = \sum_\ell v_\ell\,Z_\Omega[\ell,s]\,e_{i_\ell}$
where $e_{i_\ell}\in\R^n$ is the standard basis vector.

The composition gather$\to$scatter is:
$\operatorname{vec}(T) = \Phi^\top\Phi\operatorname{vec}(KW)$
where $\Phi\in\R^{q\times nr}$ has rows
$\Phi[\ell,:] = Z_\Omega[\ell,:]\otimes e_{i_\ell}^\top K$... more precisely,
the gather/scatter computes the action of $S^\top(Z\otimes I_n)$
followed by $(Z\otimes I_n)^\top_\Omega S$, which gives
$(Z\otimes I_n)^\top_\Omega SS^\top(Z\otimes I_n)_\Omega$ applied to
$\operatorname{vec}(KW) = (I_r\otimes K)\operatorname{vec}(W)$.

\textbf{Step 4:}
$Y = KT + \lambda KW$ applies $(I_r\otimes K)$ to $\operatorname{vec}(T)$
and adds $\lambda(I_r\otimes K)\operatorname{vec}(W)$.
The full chain: $(I_r\otimes K)\cdot\text{scatter}\cdot\text{gather}\cdot(I_r\otimes K)
+ \lambda(I_r\otimes K) = (Z\otimes K)^\top SS^\top(Z\otimes K) + \lambda(I_r\otimes K) = \mathcal{A}$.
\end{proof}

\paragraph{Comparison with direct and unpreconditioned methods.}

\begin{center}
\begin{tabular}{@{}lccc@{}}
\toprule
\textbf{Method} & \textbf{Time} & \textbf{Storage} & \textbf{$O(N)$?} \\
\midrule
Direct (form + Cholesky) & $O(n^3 r^3)$ & $O(n^2 r^2)$ & No \\
Unpreconditioned CG & $O(\kappa_\mathcal{A}(n^2 r + qr))$ & $O(nr + qr)$ & No \\
\textbf{PCG (full $K$)} & $O\bigl(\sqrt{\kappa/\rho}\ln(1/\varepsilon)(n^2r+qr+nr^2)\bigr)$ & $O(n^2+nr+qr)$ & No \\
\textbf{PCG (Nystr\"om)} & $O\bigl(\sqrt{1/\rho}\ln(1/\varepsilon)(nmr+qr+mr^2)\bigr)$ & $O(nm+nr+qr)$ & No \\
\bottomrule
\end{tabular}
\end{center}

The PCG approach achieves substantial savings over the direct method, especially
when $r$ is moderate (the $r^3$ factor in direct solve becomes $r$ or $r^2$
per iteration). The Nystr\"om variant further reduces the $n^2$ dependence
to $nm$ per iteration, which is significant for large $n$ with low-rank kernels.

\paragraph{Practical considerations.}
\begin{itemize}
  \item \textbf{Warm starting.} In the ALS outer loop, the solution from the
    previous iteration provides an excellent initial guess, reducing PCG iterations
    to typically $5$--$15$.
  \item \textbf{Landmark selection for Nystr\"om.} Uniform random selection of
    $m$ landmarks suffices for theoretical guarantees. In practice, $k$-means++
    initialization or leverage-score sampling yields better approximations.
  \item \textbf{Regularization--conditioning tradeoff.} Larger $\lambda$ improves
    the condition number (faster PCG convergence) but increases bias. In practice,
    $\lambda$ is set via cross-validation on the observed entries.
  \item \textbf{Sparse scatter.} When most mode-$k$ indices $i_\ell$ are repeated
    many times (i.e., the observation pattern is ``row-dense''), the scatter
    operation in Step~3 benefits from sorting observations by $i_\ell$, reducing
    cache misses.
\end{itemize}

\subsubsection*{Phase C: Interpretation and Verification}

\paragraph{Numerical verification.}
The complete algorithm was implemented and verified:

\begin{enumerate}
  \item \textbf{Matvec correctness.} For instances $(n, r) \in \{(8,3),(15,4),(20,5)\}$:
    the column-wise gather/scatter matvec matched the explicit $\mathcal{A}w$
    to relative error $< 10^{-15}$.
  \item \textbf{Preconditioner correctness.} The Cholesky-based Kronecker preconditioner
    $\mathcal{M}^{-1}$ matched the explicit $(G\otimes K)^{-1}$ to relative error $< 10^{-14}$.
  \item \textbf{PCG convergence.} For all test instances, PCG converged to the
    direct solution with relative error $< 10^{-11}$ and residual $< 10^{-12}$.
  \item \textbf{Condition number reduction.} The preconditioner reduced condition
    numbers substantially (e.g., from $\kappa(\mathcal{A}) \approx 10^4$ to
    $\kappa(\mathcal{M}^{-1}\mathcal{A}) \approx 50$).
  \item \textbf{Complexity scaling.} Matvec time scaled as $O(n^2 r + qr)$,
    consistent with the analysis.
\end{enumerate}

\paragraph{Confidence.} HIGH --- The algorithm rests on the classical Kronecker--vec
identity (textbook result), Conjugate Gradients for SPD systems
(Hestenes--Stiefel 1952), Cholesky-based Kronecker inversion, and the
well-studied Nystr\"om approximation. The column-wise gather/scatter
avoids $O(N)$ computation by evaluating the forward map only at observed
indices. All components were verified numerically against explicit
constructions with errors below $10^{-11}$.

%%% REFERENCES %%%

\bibitem{HestStief1952} M.~R.~Hestenes and E.~Stiefel.
Methods of conjugate gradients for solving linear systems.
\emph{J. Res. Nat. Bur. Standards}, 49(6):409--436, 1952.

\bibitem{Saad2003} Y.~Saad.
\emph{Iterative Methods for Sparse Linear Systems}.
SIAM, 2nd edition, 2003.

\bibitem{VanLoan2000} C.~F.~Van~Loan.
The ubiquitous Kronecker product.
\emph{J. Comput. Appl. Math.}, 123(1--2):85--100, 2000.

\bibitem{WilliamsSeeger2001} C.~K.~I.~Williams and M.~Seeger.
Using the Nystr\"om method to speed up kernel machines.
\emph{Advances in Neural Information Processing Systems 13}, pp.~682--688, 2001.

\bibitem{KoldaBader2009} T.~G.~Kolda and B.~W.~Bader.
Tensor decompositions and applications.
\emph{SIAM Review}, 51(3):455--500, 2009.

\bibitem{GolubVanLoan2013} G.~H.~Golub and C.~F.~Van~Loan.
\emph{Matrix Computations}.
Johns Hopkins University Press, 4th edition, 2013.

\bibitem{DrineasMahoney2005} P.~Drineas and M.~W.~Mahoney.
On the Nystr\"om method for approximating a Gram matrix for improved kernel-based learning.
\emph{J. Mach. Learn. Res.}, 6:2153--2175, 2005.

\bibitem{Acar2011} E.~Acar, D.~M.~Dunlavy, T.~G.~Kolda, and M.~M{\o}rup.
Scalable tensor factorizations for incomplete data.
\emph{Chemometrics Intell. Lab. Syst.}, 106(1):41--56, 2011.

\bibitem{Tropp2012} J.~A.~Tropp.
User-friendly tail bounds for sums of random matrices.
\emph{Found. Comput. Math.}, 12(4):389--434, 2012.

%%% HSEXPLAIN %%%

\textbf{What is this problem about?}

Imagine Netflix has a database of how users rate movies, but most of the ratings
are missing because people have only seen a small fraction of all available movies.
Netflix wants to predict what rating each user would give to each unseen movie.
The approach is to find hidden patterns: perhaps there are a few ``types'' of
movie preferences (action fans, comedy lovers, documentary enthusiasts), and each
user is a mixture of these types.

The challenge becomes harder when the patterns are not simple numbers but complex
functions. Instead of a single ``action score,'' a user's preference might depend
on subtle relationships between movie features, captured by a mathematical object
called a ``kernel.'' This kernel encodes similarity: movies that are alike in
certain ways get similar predicted ratings.

The core computational challenge is solving a massive system of equations to find
these patterns. If there are thousands of users and movies, the system could involve
millions of variables. Worse, the missing data breaks the nice mathematical structure
that would normally let us solve this quickly.

Our solution uses three ideas. First, we never write down the enormous system of
equations. Instead, we use ``matrix-free computation'': we know how to multiply by
the equation matrix without storing it, by breaking the operation into small pieces.
Each piece involves only the observed ratings (not the missing ones), using a
``gather-scatter'' technique --- gather the relevant data at observed positions,
compute a local product, then scatter the results back. The cost depends on how
many ratings were actually observed, not on the total number of possible ratings.

Second, we use a ``preconditioner'' --- a simplified version of the equations that
approximates the real system well enough that an iterative solver converges quickly.
Our preconditioner exploits a special structure called a ``Kronecker product,''
letting us work with two small matrices instead of one enormous one.

Third, we optionally compress the kernel matrix using a technique called the
Nystrom approximation. Since many kernels are effectively ``low-rank'' (most of
their information is captured by a few principal components), we can represent
them compactly, speeding up every step of the algorithm.

The result: instead of solving the equations directly (which takes time
proportional to the cube of the problem dimensions), our method takes time
roughly proportional to the number of observed ratings times the number of
patterns, with only logarithmic dependence on accuracy. This turns a computation
that might take hours into one that takes seconds.
